\section{Introduction}
\label{sec:intro}

The explosive growth of large-scale pre-trained models has led to a paradigm shift in machine learning, where a single, massive base model is adapted to a myriad of diverse, task-specific downstream applications.
To avoid the prohibitive storage and memory overhead of hosting separate model weights for every downstream client, Parameter-Efficient Fine-Tuning (PEFT) techniques, particularly Low-Rank Adaptation (LoRA) \cite{hu21lora}, have emerged as the standard industry practice.
By training only a tiny fraction of low-rank adapter parameters while keeping the large base model frozen, PEFT enables multi-tenant serving, where hundreds of specialized adapters can coexist in system storage.

However, deploying these multi-tenant systems in real-world streaming environments presents major systems challenges.
When a single server ingests a continuous, heterogeneous stream of requests from different users targeting different tasks, standard static model merging approaches---such as Task Arithmetic \cite{ilharco22arithmetic}, TIES-Merging \cite{yadav23ties}, or DARE \cite{yu23dare}---suffer from severe \emph{representation conflicts} (or ``interference'') \cite{yadav23ties}.
A single, statically merged model averaged across all tasks fails to preserve task-specific nuances, resulting in a substantial drop in task accuracy (e.g., a drop of over $15\%$ absolute accuracy compared to specialized expert models).

To resolve representation conflicts without sacrificing the memory-efficiency of merged weights, dynamic weight-merging systems have been proposed.
Among these, Parameter-Free Subspace Routing (PFSR) combined with Micro-Batch Homogenization (MBH) represents the state-of-the-art \cite{chronopoulou23representation}.
PFSR computes a dynamic, sample-specific weight interpolation coefficient on-the-fly by projecting intermediate activation features against the classification heads of specialized experts.
MBH then groups samples with similar task-routing coefficients into homogeneous micro-batches, dynamically materializing and running sequential forward passes through specialized merged models.

From a \textbf{pragmatic engineering perspective}, however, PFSR+MBH possesses a fatal systems-level flaw: a severe \textbf{two-pass latency penalty}.
Because PFSR extracts activation representations from the \emph{penultimate} layer of the network backbone, the model must execute a full, throw-away forward pass of the entire base model backbone \emph{just to compute the routing coefficients}.
Once these coefficients are determined and the downstream weights are dynamically interpolated, a \emph{second} full forward pass must be executed to obtain the actual prediction outputs.
In low-latency, real-world deployment pipelines (e.g., edge AI or high-throughput cloud servers), this doubling of inference latency ($>2.0\times$) completely negates the serving benefits of weight merging, making such penultimate-layer dynamic routing schemes highly impractical.

To bridge this critical gap between theoretical model capacity and practical systems latency, we propose \textbf{ELATI} (\textbf{E}arly-\textbf{L}ayer \textbf{A}daptive \textbf{T}ask \textbf{I}dentification).
ELATI is a training-free, parameter-free dynamic weight-merging framework designed to achieve near-optimal single-pass ($1.0\times$) latency.
We observe that the representation space of early layers (e.g., Layer 2 in a 14-layer network) contains highly structured task-specific characteristics that are already sufficient for accurate task identification.
By shifting the routing decision to an early layer $l_{\text{route}}$, ELATI eliminates the throw-away forward pass of the deep backbone.

Developing early-layer routing, however, introduces a major theoretical obstacle: because early representations have not passed through the downstream layers, there are no semantic classification heads available to compute similarity projections.
ELATI elegantly resolves this by introducing \textbf{Early-Layer Representative Mapping (ELRM)}.
Instead of projecting against class-head parameters, we perform unsupervised offline profiling on a tiny 64-sample calibration split (16 samples per task) to extract the mean activation centroids of each expert task at Layer $l_{\text{route}}$.
These centroids act as frozen, training-free projection heads, capturing the characteristic subspace manifolds of early-layer activations.

For heterogeneous streaming, we introduce \textbf{Downstream-Only Micro-Batch Homogenization (DO-MBH)}.
Incoming batches are propagated through the early shared layers $1 \dots l_{\text{route}}$ only once.
The activations at Layer $l_{\text{route}}$ are then projected against the early centroids, temperature-scaled, and partitioned into homogeneous micro-batches on-the-fly.
Finally, we dynamically interpolate the adapter weights for the downstream layers $l > l_{\text{route}}$ and dispatch each micro-batch through the merged tail of the network.
This completely avoids redundant forward computation on early layers and achieves a dramatic throughput boost.

\paragraph{Scientific Transparency regarding Experimental Setup}
We emphasize that in this work, we focus on the mathematical and systems modeling of early-layer routing.
To evaluate ELATI under controlled conditions, we implement a \textbf{Hierarchical 14-Layer Sandbox} in PyTorch.
Rather than serving actual pre-trained Vision Transformers or LLMs on GPU hardware, this sandbox mathematically structures a deep network of $L = 14$ layers with residual GeLU mappings and low-rank LoRA updates to simulate sequential layer feature transformation in a controlled, flat representation space.
We explicitly highlight that the task manifolds are modeled after the noise characteristics of standard datasets and are denoted as \emph{Simulated MNIST}, \emph{Simulated Fashion-MNIST (F-MNIST)}, \emph{Simulated CIFAR-10}, and \emph{Simulated SVHN} subspaces. This simulation setup provides a highly reproducible, mathematically transparent environment to evaluate the foundational representational dynamics of early routing.
All physical systems claims---such as VRAM scratch buffers, weight materialization, and downstream parameter interpolation---are simulated and benchmarked on CPU to analyze the relative systems complexity and complexity scaling behavior under different routing depths. On physical GPU hardware clusters, end-to-end throughput and latency are dominated by factors such as memory-bus bandwidth limits, register allocation, CUDA kernel launch overheads, and CUDA stream synchronization. Our CPU-based wall-clock timings are presented as a rigorous proof-of-concept of the computational scaling benefits and mathematical operations reductions.

Through exhaustive evaluation on this Hierarchical Sandbox across 10 independent random seeds, we demonstrate the exceptional practical utility of our proposed algorithm:
\begin{itemize}
    \item \textbf{Significant Accuracy Gains:} ELATI achieves a highly robust Joint Mean accuracy of \textbf{56.89\% $\pm$ 1.66\%}, outperforming the static Uniform Merging baseline (\textbf{48.27\% $\pm$ 2.23\%}) by \textbf{+8.62\%} absolute, demonstrating successful conflict resolution in early representation space across the simulated task activation manifolds.
    \item \textbf{Outstanding Systems Latency Benefits:} Symmetrical physical execution benchmarks on CPU demonstrate that ELATI achieves a genuine \textbf{1.40$\times$ physical CPU speedup}, reducing end-to-end execution latency of our simulated deep network from \textbf{36.90~ms} (PFSR) to \textbf{26.43~ms} on a batch of 1,000 samples, driven by avoiding 11 redundant deep base layers during Pass 1.
    \item \textbf{High Complexity Reduction:} On CPU projection step benchmarks, ELATI reduces projection latency from \textbf{1.31~ms} to \textbf{0.39~ms} under fully vectorized execution (a \textbf{3.33$\times$ speedup}) and from \textbf{171.05~ms} to \textbf{63.48~ms} under a sequential loop setup (a \textbf{2.69$\times$ speedup}). We present a symmetric evaluation of these paradigms, and explain how the persistent speedup is driven by a reduction in mathematical complexity from $O(B \cdot K \cdot C \cdot D)$ to $O(B \cdot K \cdot D)$.
    \item \textbf{High-Scarcity Robustness:} Unlike classical Linear Routers which overfit and collapse under data scarcity, ELATI's unsupervised similarity-based centroid projection maintains exceptional stability across all simulated tasks with zero parameter optimization and is highly robust under representational entanglement sweeps.
\end{itemize}

By offering an optimized trade-off between task performance and actual systems complexity, ELATI establishes a highly practical, deployable path for multi-tenant dynamic weight merging.
